% Input:       nothing; this file holds no code of its own.
% Output:      the six kernel stages, input in pipeline order.
% Owned state: none.
% Invariants:  the stage order is the pipeline order; each stage guards itself.
% Next stage:  tenkz-kernel-language.code.tex, the first stage listed below.
% tenkz-kernel.code.tex
%
% The 1.0 kernel loader.  This file held the whole kernel until the stage
% split; it now names the six stages and their order, and nothing else.  The
% order is the pipeline: the language stage validates records, the policy
% stage derives the wires a picture's own words imply, resolution answers
% every spatial question, ink lays the marks down, annotation decides where a
% name or a mark may stand, and the surface stage exposes the environments.
%
% Consumers that input the kernel by name -- the kernel and regression
% fixtures, and the API examples -- keep working unchanged, which is why the
% name stays.  Each stage carries its own load guard.
%
% SPDX-License-Identifier: Apache-2.0
% Copyright the TNLean project; see LICENSE for the full terms.

% Load guard, as the single file carried: a document that inputs the kernel
% a second time must not walk the six stages again.
\ExplSyntaxOn
\cs_if_exist:NT \tenkzkernel { \tex_endinput:D }
\ExplSyntaxOff

% Each stage closes its own \ExplSyntaxOn, so the syntax is re-armed before
% every input rather than assumed to survive the one before it.
\ExplSyntaxOn \file_input:n { tenkz-kernel-language.code.tex }
\ExplSyntaxOn \file_input:n { tenkz-kernel-policy.code.tex }
\ExplSyntaxOn \file_input:n { tenkz-kernel-resolve.code.tex }
\ExplSyntaxOn \file_input:n { tenkz-kernel-ink.code.tex }
\ExplSyntaxOn \file_input:n { tenkz-kernel-annotate.code.tex }
\ExplSyntaxOn \file_input:n { tenkz-kernel-surface.code.tex }
\endinput
